\section{Proof Details}
\label{sec:proofs}

\subsection{Subspace geometry: cross-Gram, budget identity, and collapse}
\label{sec:geometry-proofs}

\paragraph{Closed-form cross-Gram.} For $\Delta\sim\mathrm{Unif}[0,L]$ and
$x_\omega(\Delta)=[\cos(\omega\Delta)\ \ \sin(\omega\Delta)]$, write
$d=(\omega{-}\nu)L$, $s=(\omega{+}\nu)L$, $a(t)=\sin t/t$ and
$b_\star(t)=(1{-}\cos t)/t$. Expanding the products of sinusoids and integrating
termwise gives
\begin{equation}
H_{\omega\nu}=\mathbb E[x_\omega^\top x_\nu]=\frac12
\begin{bmatrix}
a(d)+a(s) & b_\star(s)-b_\star(d)\\
b_\star(s)+b_\star(d) & a(d)-a(s)
\end{bmatrix},
\label{eq:cross-gram}
\end{equation}
and $S_\omega=H_{\omega\omega}$. Since $Q_{\omega\nu}=S_\omega^{-1/2}H_{\omega\nu}S_\nu^{-1/2}$
is the whitened cross-Gram of two two-dimensional subspaces, its singular values
are the canonical correlations, and $c_{\omega\nu}=\frac12\|Q_{\omega\nu}\|_F^2$
is invariant under $x_\omega\mapsto x_\omega O$ for any invertible $O$ (in
particular under rotation of the phase inside a pair), because $S_\omega$
transforms congruently and the whitening cancels $O$.

\begin{proof}[Proof of Theorem~\ref{thm:budget}]
$R$ is the block matrix with blocks $R_{ij}=Q_{\omega_i\omega_j}$ and
$R_{ii}=I_2$. Hence $\operatorname{tr}R=\sum_{i=1}^K\operatorname{tr}I_2=2K$.
For the second moment,
$\operatorname{tr}(R^2)=\sum_{i,j}\operatorname{tr}(R_{ij}R_{ji})
=\sum_i\operatorname{tr}(I_2)+\sum_{i\neq j}\|Q_{\omega_i\omega_j}\|_F^2
=2K+2K(K-1)\bar c$, using $R_{ji}=R_{ij}^\top$ and the definition
$\bar c=\frac{1}{K(K-1)}\sum_{i\neq j}c_{\omega_i\omega_j}$ with
$\|Q_{ij}\|_F^2=2c_{ij}$. Dividing gives \eqref{eq:budget-identity}.
\end{proof}

The identity is exact and involves only the mean of the pairwise redundancies;
it does not claim that $\bar c$ determines the Shannon effective rank or the
log-determinant, which retain higher-order multi-subspace dependence.

\begin{proof}[Proof of Proposition~\ref{prop:collapse}]
With $t=\Delta/L$ and $x=\omega L$, $\cos(xt)=1-x^2t^2/2+O(x^4)$ and
$\sin(xt)/x=t-x^2t^3/6+O(x^4)$, so after rescaling the second coordinate the
basis of $V_\omega$ tends to $\{1,\Delta\}$. Substituting these expansions into
\eqref{eq:cross-gram} and whitening, the leading deficit between two slow
subspaces is $2-\|Q_{x,y}\|_F^2=\frac{19}{12600}(x^2-y^2)^2+O(\epsilon^6)$;
numerically at $x{=}0.05,y{=}0.10$ the exact-to-leading ratio is $1.00058$. For
the softmax metric, $F=\operatorname{diag}(p)-pp^\top$ satisfies $F\mathbf1=0$,
so constants are annihilated; centring by $\mathbb E_p$ and rescaling gives
$\overline{\sin(\omega\Delta)}/\omega\to\Delta-\mathbb E_p\Delta$ and
$-2\overline{\cos(\omega\Delta)}/\omega^2\to\Delta^2-\mathbb E_p\Delta^2$,
which span the centred limit whenever $p$ has nondegenerate support on at least
three distances.
\end{proof}

The reported $r_2{=}2.00$ for the $23$ slow pairs uses the standard
$u_k{=}k/K$ grid and the uniform separation
prior above. Other priors change that static numerical Gram; the fixed-range
trained-model contrast of \S\ref{sec:exp-identify} does not use this prior.

\subsection{Post-hoc transplant obstruction}
\label{sec:obstruction-proof}

\begin{proof}[Proof of Theorem~\ref{thm:obstruction}]
Both $R_\Omega$ and $R_{\Omega'}$ are block-diagonal with $2\times2$ rotation
blocks, hence $R_\Omega(0)=R_{\Omega'}(0)=I$. Evaluating the hypothesis at
$\Delta{=}0$ gives $A^\top B=I$, so $B=A^{-\top}$ and the hypothesis becomes the
similarity $A^\top R_{\Omega'}(\Delta)A^{-\top}=R_\Omega(\Delta)$ on an interval.
Each side is the exponential of a constant generator, $R_\Omega(\Delta)=\exp(\Delta G_\Omega)$
with $G_\Omega=\bigoplus_k \omega_k J$, $J=\left[\begin{smallmatrix}0&-1\\1&0\end{smallmatrix}\right]$.
Differentiating at $\Delta{=}0$ gives $A^\top G_{\Omega'}A^{-\top}=G_\Omega$, so
the generators are similar and share a spectrum. The spectrum of $G_\Omega$ is
$\{\pm i\omega_k\}_{k\le K}$ with multiplicity, so the frequency multisets agree
up to sign and permutation. When a frequency is repeated, similarity may mix
the entire equal-frequency invariant subspace; it need not preserve the original
two-dimensional block decomposition. For integer positions the same argument
is replaced by a one-step spectral argument: evaluating at $\Delta=0$ again
gives $B=A^{-\top}$, while $\Delta=1$ makes $R_{\Omega'}(1)$ and
$R_\Omega(1)$ similar. Their eigenvalue multisets
$\{e^{\pm i\omega_k}\}$ therefore agree, which identifies frequencies up to
sign, permutation, and the $2\pi$ alias.
\end{proof}

This is an exact statement about \emph{fixed, position-independent, invertible}
compensation. It does not bound approximate compensation obtained by retraining
on finite data, nor does it constrain nonlinear maps or the V/O and residual
paths.

\subsection{Collision-only surrogate stationary density}

This subsection proves Theorem~\ref{thm:ode} directly from the stated
collision-only surrogate functional. Starting from the formal surrogate kernel
\begin{equation}
K_{\mathrm{app}}(\phi,\psi)=\alpha\delta(\phi-\psi)+\beta\min(\phi,\psi),
\end{equation}
the main-text collision functional is~\eqref{eq:Capp}, which under the mass normalization $\int_0^1\rho=1$ with Lagrange multiplier $\nu$ has constrained first variation
\begin{equation}
\alpha\rho(\phi) + \beta\,g(\phi) + \nu = 0,\qquad
g(\phi) = \int_0^1 \rho(\psi)\min(\phi,\psi)\,d\psi.
\end{equation}
\paragraph{Existence.}
For rigor, view $\mathcal{C}_{\mathrm{app}}$ first as a functional on $L^2([0,1])$ and minimize over $\mathcal{A}=\{\rho\in L^2([0,1]):\rho\ge0\ {\rm a.e.},\int_0^1\rho=1\}$. The $\alpha\|\rho\|_2^2$ term is coercive and weakly lower semicontinuous, the Green-kernel term is positive semidefinite and continuous, and the mass/nonnegativity constraints are weakly closed. A minimizer therefore exists by the direct method. Strict convexity gives uniqueness. The Euler solution derived below is strictly positive, so the nonnegativity constraint is inactive a posteriori and the classical $C^2$ derivation is justified.

Direct differentiation of $g$ gives
\begin{align}
g'(\phi)
&=
\int_\phi^1 \rho(\psi)\,d\psi,
&
g''(\phi)
&=
-\rho(\phi).
\end{align}
Differentiating the first variation twice in $\phi$ eliminates $\nu$ (a constant) and yields the homogeneous ODE
\begin{equation}
\rho''(\phi) - \tau^2\rho(\phi) = 0,\qquad \tau = \sqrt{\beta/\alpha}.
\label{eq:homogeneous-ode}
\end{equation}
Evaluating the pre-differentiated stationarity together with $g(0){=}0$, $g'(1){=}0$, and $\int_0^1\rho = 1$ gives the derivative boundary conditions
\begin{equation}
\rho'(0) = -\tau^2,\qquad \rho'(1) = 0,
\label{eq:bcs}
\end{equation}
since differentiating $\alpha\rho(\phi)+\beta g(\phi)+\nu = 0$ once and evaluating at $\phi{=}1$ gives $\alpha\rho'(1)+\beta g'(1) = 0$ hence $\rho'(1){=}0$, and at $\phi{=}0$ gives $\alpha\rho'(0)+\beta g'(0) = 0$ with $g'(0) = \int_0^1\rho = 1$, so $\rho'(0) = -\beta/\alpha = -\tau^2$. Substituting $\rho(\phi) = C_1\cosh(\tau\phi)+C_2\sinh(\tau\phi)$ into~\eqref{eq:homogeneous-ode} and~\eqref{eq:bcs} together with $\int_0^1\rho = 1$ gives the unique positive interior stationary density
\begin{equation}
\rho_\tau(\phi) \;=\; \frac{\tau\,\cosh(\tau(1-\phi))}{\sinh\tau},
\label{eq:rho-tau-closed}
\end{equation}
which is main-text Theorem~\ref{thm:ode}. The mass constraint $\int\rho{=}1$ is not lost: it enters the boundary condition at $\phi{=}0$ via $g'(0){=}\int_0^1\rho{=}1$. Positivity holds on the closed interval since $\rho_\tau(\phi){\geq}\rho_\tau(1){=}\tau/\sinh\tau{>}0$ for $\tau{>}0$ (with $\rho_\tau(0){=}\tau\coth\tau{>}0$); the inequality constraint $\rho{>}0$ is therefore inactive throughout $[0,1]$ and the unconstrained Euler--Lagrange solution coincides with the KKT-constrained solution. For the degenerate case $\beta{=}0$ the ODE reduces to $\rho''{=}0$ with $\rho'(0){=}\rho'(1){=}0$ and $\int\rho{=}1$, forcing $\rho\equiv 1$, which is also the $\tau{\to}0$ limit of~\eqref{eq:rho-tau-closed}.

\paragraph{Convexity / PSD of $\mathcal{C}_{\mathrm{app}}$.} For $\alpha{>}0$ and $\beta{\geq}0$, the surrogate is convex because both quadratic terms are positive semidefinite. The diagonal $\alpha{\int}\rho^2$ is manifestly PSD on $L^2$, and the Green-kernel cross term satisfies the identity
\begin{equation}
\iint_{[0,1]^2} f(\phi)\,f(\psi)\,\min(\phi,\psi)\,d\phi\,d\psi \;=\; \int_0^1 \!\left(\int_s^1 f(u)\,du\right)^{\!2}\!ds \;\geq\; 0.
\label{eq:min-kernel-psd}
\end{equation}
Thus the stationary density obtained above is the unique constrained \emph{minimizer} (not merely a critical point), and the surrogate is well-defined as a quadratic functional on $L^2([0,1])\supset C^2([0,1])$ with no further regularity assumption needed for the $\delta$-component. Conditional on $\mathcal{C}_{\mathrm{app}}$, this derivation is exact.

\paragraph{Why this branch.}
The main construction uses the homogeneous Cosh solution because its CDF is
analytically invertible. Adding a task-dependent forcing term would produce a
different, generally numerical allocation. We therefore test the closed-form
branch functionally below rather than treating it as an approximation to an
unknown trained-loss optimum.

\subsection{Surrogate quality: functional validation}
\label{sec:surrogate-validation}

The broadband surrogate $K_{\mathrm{app}} = \alpha\delta + \beta\min$ is the
primary approximation behind the closed-form density; selecting the pure-tether
branch and a finite operating point are additional modeling choices. We test the
surrogate functionally rather than pointwise against the closed-form cosine-feature kernel
$K_{\cos}(\phi_1,\phi_2) = \int D(\Delta)\cos(\omega(\phi_1)\Delta)\cos(\omega(\phi_2)\Delta)\,d\Delta$:
does the derived allocation also reduce the collision objective under $K$? This
checks the direction selected by the homogeneous branch; it does not bound the
omitted Fisher-forcing residual.

Table~\ref{tab:surrogate-validation} answers affirmatively. For a uniform distance prior $D(\Delta) = 1/L$ on $[0, L]$, we compute $K_{\cos}$ at the discrete allocation points for both geometric and EVQ, then evaluate the normalized collision score $C = \sum_{i<j} K_{ij}^2/(K_{ii}K_{jj})$ and effective spectral rank (exponential of the Shannon entropy of the eigenvalue distribution). These diagnostics use the inclusive grid $u_k=k/(K-1)$; they test the allocation direction but are not a direct validation of either the full block-whitened sin/cos objective or the deployed midpoint grid. The $12$ reported configurations comprise a six-length text sweep at fixed $b{=}500$K and $K{=}32$, two additional text bases at $L{=}2048$, and four video bases at $L{=}32$ and $K{=}16$; this is not a Cartesian sweep over every listed factor.

\begin{table}[ht]
\caption{Functional surrogate check on the inclusive $u_k{=}k/(K{-}1)$ diagnostic grid. EVQ reduces collision and increases effective rank under the closed-form cosine-feature kernel $K_{\cos}$ in all 12 tested configurations. This is a direction check, not a pointwise approximation, full-subspace validation, or deployed-grid identity test.}
\label{tab:surrogate-validation}
\centering
\small
\begin{tabular}{@{}l r r r r r@{}}
\toprule
& \multicolumn{2}{c}{Collision $C$} & & \multicolumn{2}{c}{Eff.\ rank} \\
\cmidrule(lr){2-3} \cmidrule(lr){5-6}
Configuration & Geo & \evq{} ($\Delta_C$) & & Geo & \evq{} \\
\midrule
\multicolumn{6}{@{}l}{\emph{Text: $K{=}32$, $b{=}500\mathrm{K}$, sweep $L$}} \\
\quad $L{=}128$ & 226.3 & 17.7 ($-92\%$) & & 2.5 & 17.0 \\
\quad $L{=}256$ & 192.6 & 27.6 ($-86\%$) & & 3.0 & 13.9 \\
\quad $L{=}512$ & 163.2 & 43.3 ($-73\%$) & & 3.6 & 10.8 \\
\quad $L{=}1024$ & 133.9 & 59.0 ($-56\%$) & & 4.4 & 8.8 \\
\quad $L{=}2048$ & 109.0 & 67.5 ($-38\%$) & & 5.4 & 8.0 \\
\quad $L{=}4096$ & 87.0 & 66.3 ($-24\%$) & & 6.6 & 8.1 \\
\midrule
\multicolumn{6}{@{}l}{\emph{Text: $K{=}32$, $L{=}2048$, sweep $b$}} \\
\quad $b{=}10\mathrm{K}$ & 34.8 & 19.2 ($-45\%$) & & 13.0 & 17.3 \\
\quad $b{=}100\mathrm{K}$ & 79.4 & 47.0 ($-41\%$) & & 7.1 & 10.4 \\
\midrule
\multicolumn{6}{@{}l}{\emph{Video: $K{=}16$, $L{=}32$, sweep $b$}} \\
\quad $b{=}100$ & 22.6 & 14.2 ($-37\%$) & & 4.8 & 6.4 \\
\quad $b{=}1\mathrm{K}$ & 42.3 & 28.0 ($-34\%$) & & 2.8 & 3.9 \\
\quad $b{=}10\mathrm{K}$ & 55.4 & 39.8 ($-28\%$) & & 2.2 & 2.9 \\
\quad $b{=}50\mathrm{K}$ & 63.3 & 47.2 ($-26\%$) & & 2.0 & 2.5 \\
\bottomrule
\end{tabular}
\end{table}

EVQ reduces collision under $K_{\cos}$ in all 12 configurations ($24\%$--$92\%$), and effective rank increases in every row. In the fixed-base text sweep, the relative reduction grows as $L$ shortens. The separate base sweeps remain positive but are not monotone in $b$.

\subsection{Single-crossing budget shift}
\label{sec:budget-crossing-proof}

\begin{lemma}[Single-crossing budget shift]
\label{lem:budget-crossing}
For every $\tau>0$, $\rho_\tau$ crosses the uniform density exactly once, at
\begin{equation}
\phi_c(\tau)=1-\tau^{-1}\operatorname{arcosh}(\sinh\tau/\tau)
\le 1-1/\sqrt3.
\label{eq:budget-crossing}
\end{equation}
\end{lemma}

\begin{proof}[Proof of Lemma~\ref{lem:budget-crossing}]
For $\tau>0$, $\rho_\tau(\phi)=\tau\cosh(\tau(1-\phi))/\sinh\tau$ is
strictly decreasing on $[0,1)$, with
$\rho_\tau(0)=\tau\coth\tau>1$ and
$\rho_\tau(1)=\tau/\sinh\tau<1$.  Hence the crossing is unique, and solving
$\rho_\tau(\phi_c)=1$ gives~\eqref{eq:budget-crossing}.  To bound it, expand
\[
\frac{\sinh\tau}{\tau}
=\sum_{n\ge0}\frac{\tau^{2n}}{(2n+1)!}
\ge \sum_{n\ge0}\frac{\tau^{2n}}{3^n(2n)!}
=\cosh\!\left(\frac{\tau}{\sqrt3}\right),
\]
where the coefficient-wise inequality is $3^n\ge2n+1$.  Monotonicity of
$\operatorname{arcosh}$ yields
$\operatorname{arcosh}(\sinh\tau/\tau)\ge\tau/\sqrt3$, proving the bound.
The three-region masses follow directly from
$F_\tau(\phi)=1-\sinh(\tau(1-\phi))/\sinh\tau$ at
$\phi_L=\log(L/2\pi)/\log b$ and $\phi_M=\log(M/2\pi)/\log b$.
\end{proof}

\subsection{Waterbed envelope and corollary}
\label{sec:waterbed-proof}

\paragraph{Setup.} With $\rho{>}0$, $\int_0^1\rho{=}1$, $U(d\phi){=}d\phi$, $P_\rho(d\phi){=}\rho\,d\phi$, $\delta\rho{=}\rho{-}1$, and $M_\rho(s){=}\int_s^1\delta\rho\,du$, define the waterbed functional
\begin{equation}
\mathcal{W}(\rho) \;\coloneqq\; \alpha\,\chi^2(U\|P_\rho) + \beta\int_0^1 M_\rho(s)^2\,ds,\qquad \chi^2(U\|P_\rho) = \int_0^1 \frac{(1-\rho)^2}{\rho}\,d\phi.
\label{eq:waterbed-functional}
\end{equation}
The first term is the Pearson divergence of $U$ from $P_\rho$ (equivalently, $\int 1/\rho - 1$); the second is the integrated squared cumulative error, which arises from the $\min$-kernel off-diagonal covariance in $K_{\mathrm{app}}$.

\begin{proposition}[Waterbed inequality]
\label{prop:waterbed}
Let $D_B(\rho) = -\int_0^1\log\rho(\phi)\,d\phi$ be the Burg entropy deficit. Then
\begin{equation}
\mathcal{W}(\rho) \;\ge\; \alpha\,F_{*}(D_B(\rho)),\qquad F_{*}(d) \;=\; \inf_{\rho>0,\,\int\rho=1,\,-\int\log\rho=d}\left[\int_0^1\frac{1}{\rho}-1\right],
\label{eq:waterbed-tight}
\end{equation}
with equality $\mathcal{W}(\rho) = 0$ if and only if $\rho \equiv 1$ (geometric RoPE). A readable corollary via R\'enyi monotonicity is
\begin{equation}
\mathcal{W}(\rho) \;\ge\; \alpha\bigl(e^{D_B(\rho)} - 1\bigr).
\label{eq:waterbed-corollary}
\end{equation}
\end{proposition}

\begin{proof}[Proof sketch]
Since $\beta\int M_\rho^2\ge 0$, $\mathcal{W}(\rho)\ge\alpha\chi^2(U\|P_\rho)$. For fixed $D_B(\rho){=}d$, the sharp lower envelope of $\chi^2$ over all densities is the Csisz\'ar--Vajda $f$-divergence comparison $F_{*}$, whose extremizer is a two-level density $X=1/\rho\in\{a,c\}$ with mass-normalization and entropy-deficit constraints. For the corollary, $1+\chi^2(U\|P_\rho) = \int 1/\rho\,\mathbb{E}_{U}[\cdot] \ge e^{D_1(U\|P_\rho)} = e^{D_B(\rho)}$ by R\'enyi-2 $\ge$ R\'enyi-1. Equality at $\rho{\equiv}1$ is immediate from both sides vanishing.
\end{proof}

\paragraph{EVQ evaluation.} For $\rho_\tau(\phi) = \tau\cosh(\tau(1{-}\phi))/\sinh\tau$ (Eq.~\eqref{eq:S-chi2-evq} below uses the un-normalized integral; cf.\ \S\ref{sec:chi2-load} ``Normalization convention''),
\begin{align}
D_B(\rho_\tau) \;&=\; \log\tfrac{\sinh\tau}{\tau} - \tfrac{1}{\tau}\int_0^\tau\log\cosh y\,dy \;=\; \tfrac{\tau^4}{90} - \tfrac{8\tau^6}{2835} + O(\tau^8),\label{eq:D-B-evq}\\
S_{\chi^2}(\rho_\tau) \;&=\; \tfrac{\sinh\tau}{\tau^2}\arctan(\sinh\tau) - 1 \;=\; \tfrac{\tau^4}{45} - \tfrac{2\tau^6}{315} + O(\tau^8),\label{eq:S-chi2-evq}\\
\int_0^1 M_{\rho_\tau}(s)^2\,ds \;&=\; \tfrac{2\tau^4}{945} + O(\tau^6).\label{eq:M-rho-evq}
\end{align}
Combining,
\begin{equation}
\mathcal{W}(\rho_\tau) \;=\; \left(\frac{\alpha}{45}+\frac{2\beta}{945}\right)\tau^4 + O(\tau^6),\qquad \mathcal{W}(\rho_{\tau_*}) \;=\; O\!\left(\frac{d_{\mathrm{head}}^4}{(L_{\mathrm{eff}}^{J})^2}\right).
\label{eq:waterbed-evq-scale}
\end{equation}
The centered companion functional is $O(\tau^4)$. It is not the value difference of $\mathcal{C}_{\mathrm{app}}$ itself and does not determine trained-model NLL.

\subsection{Surrogate self-consistency theorem}
\label{sec:self-consistency}

The surrogate $\mathcal{C}_{\mathrm{app}}$ contains two integral terms $T_1(\tau) = \int_0^1 \rho^2\,d\phi$ and $T_2(\tau) = \int_0^1\!\!\int_0^1 \rho(\phi)\rho(\psi)\min(\phi,\psi)\,d\phi\,d\psi$, where $\rho(\phi) = \tau\cosh(\tau(1{-}\phi))/\sinh\tau$ is the normalized cosh density.

\begin{theorem}[Surrogate self-consistency]
\label{thm:self-consistency}
For all $\tau > 0$:
\begin{equation}
\tau^2 T_2(\tau) + T_1(\tau) = \tau\coth\tau.
\label{eq:self-consistency}
\end{equation}
\end{theorem}

\begin{proof}
The Green's function $g(\phi) = \int_0^1 \rho(\psi)\min(\phi,\psi)\,d\psi$ satisfies $g''(\phi) = -\rho(\phi)$ with $g(0) = 0$, $g'(1) = 0$ (from the $\min$ kernel structure; see~\S\ref{sec:proofs}). Define $h(\phi) = \rho(0) - \rho(\phi)$. Then $h'' = -\rho''$, and under the pure-tether ODE $\rho'' = \tau^2\rho$ this gives $h'' = -\tau^2\rho$, with $h(0) = 0$ and $h'(0) = -\rho'(0)$.

Define $f(\phi) = \tau^2 g(\phi) - h(\phi)$. Then $f'' = \tau^2 g'' - h'' = -\tau^2\rho - (-\tau^2\rho) = 0$, so $f$ is linear in $\phi$. The boundary values give $f(0) = \tau^2 g(0) - h(0) = 0$ and $f'(0) = \tau^2 g'(0) + \rho'(0) = \tau^2 - \tau^2 = 0$ (using $g'(0) = \int_0^1\rho = 1$ and $\rho'(0) = -\tau^2$ from the pure-tether density). Hence $f \equiv 0$, yielding $\tau^2 g(\phi) = h(\phi) = \rho(0) - \rho(\phi)$, i.e.,
\begin{equation}
g(\phi) = \frac{\rho(0) - \rho(\phi)}{\tau^2}.
\label{eq:green-identity}
\end{equation}
Integrating against $\rho$:
\begin{equation}
T_2 = \int_0^1 \rho(\phi)\,g(\phi)\,d\phi = \frac{1}{\tau^2}\bigl[\rho(0)\underbrace{\int_0^1\!\rho}_{=1} - T_1\bigr]
= \frac{\rho(0) - T_1}{\tau^2}.
\end{equation}
Since $\rho(0) = \tau\cosh\tau/\sinh\tau = \tau\coth\tau$, rearranging gives~\eqref{eq:self-consistency}.
\end{proof}

An immediate corollary follows.

\begin{corollary}[Closed-form $T_2$]
\begin{equation}
T_2(\tau) = \frac{\sinh 2\tau - 2\tau}{4\tau\sinh^2\tau}.
% \label{eq:T2-closed} % unused
\end{equation}
\end{corollary}

This follows by substituting
$T_1(\tau)=\tau^2(1+\sinh(2\tau)/(2\tau))/(2\sinh^2\tau)$ and
$\rho(0)=\tau\coth\tau$ into the theorem.

The identity is an exact internal check on the two surrogate energy terms; it makes no trained-loss claim.

\subsection{From scaling structure to an operating default}
\label{sec:lambda-cv}

The post-softmax model uses
$\mathcal{F}(\tau)=\tfrac12S_{\chi^2}(\tau)-\lambda U(\tau,L)$ with
$S_{\chi^2}(\tau)=\tau^4/(45d_{\mathrm{head}})+O(\tau^6)$ and
$U(\tau,L)=(d_{\mathrm{head}}/L)[Q_0+\tau^2Q_1(L,b)+O(\tau^4)]$.
Its leading stationary balance is
\begin{equation}
\tau_*^2=45\lambda Q_1(L,b)\frac{d_{\mathrm{head}}^2}{L}.
\label{eq:tau-leading-balance}
\end{equation}
Thus the model supplies the powers of $d_{\mathrm{head}}$ and $L$; the
finite coefficient $c(\Pi)=\sqrt{45\lambda Q_1}$ depends on the chosen
utility, base, and protocol. Setting $c{=}1$ is the deployed zero-search
convention, not an independently derived constant.

Here
\[
q(x)=\tfrac12+\frac{\sin(2x)}{4x}-\left(\frac{\sin x}{x}\right)^2
=\operatorname{Var}_{t\sim U[0,1]}[\cos(xt)]
\]
and $Q_1(L,b)=\int_0^1\eta(\phi)q(Lb^{-\phi})\,d\phi$. This is a
specified phase-variance utility. It should not be identified with the exact
oscillatory collision objective or the trained-transformer loss.

The finite-$\tau$ evidence is deliberately empirical. A staged 99-run study
contains nine small-model configurations and 27 formula-versus-Geo seed pairs,
not 27 independent configurations. The default beats midpoint-Geo in $7/9$
configuration means and $18/27$ pairs, but beats a pilot-selected neighbouring
value in only $3/9$ held-out configuration means. In the separate exact-range
factorial, the formula point is the best pre-specified Cosh multiplier in
$4/12$ structural configurations. These audits establish the correct role of
$c{=}1$: a useful, fallible basin initialisation.

\subsection{Scaling analysis of \texorpdfstring{$\tau^*$}{tau*}: from surrogate to operating point}
\label{sec:tau-scaling}

\paragraph{Numerical characterization of surrogate coefficients.}
We fit $K_{\mathrm{app}}(\phi_i,\phi_j) = \alpha\delta_{ij}/\Delta\phi + \beta\min(\phi_i,\phi_j)$ to the exact collision kernel $K(\phi_i,\phi_j) = \frac{1}{2L}\bigl[\frac{\sin((\omega_1{-}\omega_2)L)}{\omega_1{-}\omega_2} + \frac{\sin((\omega_1{+}\omega_2)L)}{\omega_1{+}\omega_2}\bigr]$ evaluated at $K = d_{\mathrm{head}}/2$ uniformly spaced channel positions in $[0,1]$. Sweeping across $L \in \{128,\ldots,4096\}$, $d_{\mathrm{head}} \in \{32,64,128,256\}$, and $b \in \{10\mathrm{K}, 100\mathrm{K}, 500\mathrm{K}\}$ yields:
\begin{equation}
\alpha \;\approx\; \frac{1}{2K} \;=\; \frac{1}{d_{\mathrm{head}}},
\qquad
\beta \;\sim\; L^{-0.22},\; O(1).
\end{equation}
The diagonal coefficient $\alpha$ is set by the kernel's diagonal value ($K(\phi,\phi) \approx 1/2$ for most channels) times the channel spacing $\Delta\phi = 1/K$. Crucially, $\alpha$ depends on $d_{\mathrm{head}}$ through the discretization ($K = d_{\mathrm{head}}/2$) and is nearly independent of $L$. The off-diagonal coefficient $\beta$ has weak power-law decay with $L$ and negligible $d_{\mathrm{head}}$ dependence.

\paragraph{The two parameters called $\tau$.}
Substituting into $\tau_{\mathrm{surr}} = \sqrt{\beta/\alpha}$:
\begin{equation}
\tau_{\mathrm{surr}} = \sqrt{\beta \cdot d_{\mathrm{head}}} \;\sim\; \sqrt{d_{\mathrm{head}}} \cdot L^{-0.11}.
\end{equation}
This fitted-ratio scaling is neither the deployed rule nor evidence for it. The
broadband surrogate selects the Cosh family; the separate post-softmax model in
Proposition~\ref{prop:softmax-transport} balances an $O(\tau^2)$ utility with
an $O(\tau^4)$ stiffness and yields
\begin{equation}
\tau_*^2 \propto \frac{d_{\mathrm{eff}}^2}{L}.
\end{equation}
The equality $c{=}1$ in the deployed rule is then evaluated by trained-model
controls as described in \S\ref{sec:lambda-cv}; it is not obtained from the
fitted $\beta/\alpha$ ratio.

\subsection{Post-softmax operating-point model}
\label{sec:softmax-transport}

The deployed scale is selected by a model separate from the surrogate that
selects the shape. We state it here rather than in the main text because its
small-$\tau$ expansion does not certify the deployed finite $\tau$.

\begin{proposition}[Semi-analytic operating-point scaling]
\label{prop:softmax-transport}
On the pure-tether family, assuming $Q_1(L,b)>0$ and a diffuse post-softmax
baseline $p_0{=}1/L$, the Pearson-$\chi^2$ channel-load stiffness and the
post-softmax transport utility admit small-$\tau$ expansions
$S_{\chi^2}(\tau) = c_S\tau^4/\dhd + O(\tau^6)$ with $c_S{=}\tfrac{1}{45}$, and
$U(\tau,L) = (\dhd/L)[Q_0 + \tau^2 Q_1(L,b) + O(\tau^4)]$, where
$Q_1(L,b) = \int_0^1 \eta(\phi)q(Lb^{-\phi})d\phi$,
$\eta(\phi) = (1{-}\phi)^2/2 - 1/6$ and
$q(x) = 1/2 + \sin(2x)/(4x) - (\sin x/x)^2$. The leading-order stationary point
of $\mathcal{F}(\tau) = \tfrac12 S_{\chi^2}(\tau) - \lambda U(\tau,L)$ satisfies
\[
\tau_*^2 = 45\lambda Q_1(L,b)\cdot\frac{\dhd^2}{L},
\qquad \tau_* \propto \frac{\dhd}{\sqrt{L}}.
\]
\end{proposition}

The explicit $\dhd/\sqrt{L}$ factor is modulated by $Q_1(L,b)$, so this is a
scaling structure rather than a pure power law, and the prefactor
$c(\Pi){=}\sqrt{45\lambda Q_1(L,b)}$ remains convention- and
protocol-dependent. Two caveats bound what the proposition supports. The
expansion is a small-$\tau$ statement, whereas the deployed configurations use
$\tau$ near $4$; at those values the exact cosh warp is evaluated directly and
the series only motivates the functional form. The diffuse baseline $p_0{=}1/L$
is also a modelling assumption: trained attention is neither uniform nor
stationary across heads, and the softmax-Jacobian effective window
$L_{\mathrm{eff}}^J$ of \S\ref{sec:leff-derivation} is the object that would
replace it, but we do not measure it here. The deployed $c{=}1$ is therefore an
empirical operating default, evaluated by the trained-model controls of
\S\ref{sec:exp-tau}, not a consequence of this proposition.

\subsection{Stiffness functional derivation}
\label{sec:stiffness-derivation}

The softmax transport objective in Proposition~\ref{prop:softmax-transport} requires choosing a stiffness functional $S(\tau)$ measuring the cost of departing from uniform density. We test a one-parameter family $S_p(\tau) = \frac{1}{d_{\mathrm{head}}}\int_0^1 (\rho_\tau - 1)^2/\rho_\tau^p\,d\phi$ using exact (non-Taylor) integrals of the normalized cosh density $\rho_\tau(\phi) = \tau\cosh(\tau(1{-}\phi))/\sinh\tau$.

\begin{table}[ht]
\caption{$L$-exponent $\gamma$ (defined by $\tau_* \propto L^{-\gamma}$) under different stiffness functionals, computed via exact numerical optimization over $L \in [128, 4096]$. Target: $\gamma = 0.500$.}
\label{tab:stiffness-sweep}
\centering
\small
\begin{tabular}{@{}l c c l@{}}
\toprule
Stiffness & $p$ & $\gamma$ & Note \\
\midrule
$L^2$ norm $\int(\rho{-}1)^2$ & 0 & 0.626 & Original choice \\
Geometric mean & 0.5 & 0.561 & \\
Exponent-matched diagnostic & 0.80 & 0.498 & Sensitivity analysis, not used in deployment \\
\textbf{Pearson $\chi^2$} & \textbf{1.0} & \textbf{0.465} & Channel-load variance penalty (axioms) \\
\midrule
KL divergence & --- & 0.710 & Worse than $L^2$ \\
Reverse KL & --- & 0.581 & \\
Jeffreys & --- & 0.621 & $\approx L^2$ \\
R\'enyi ($\alpha{=}2$) & --- & 0.807 & \\
\bottomrule
\end{tabular}
\end{table}

\paragraph{Why $\chi^2$.} Two independent arguments favor $p{=}1$. (1)~\emph{Attention-load asymmetry}: each frequency channel at density $\rho(\phi)$ bears per-channel attention load ${\propto}1/\rho$; weighting the squared departure by this load yields $\int(\rho{-}1)^2/\rho = \chi^2$. (2)~\emph{Discrete-grid effect}: the continuous $L^2$ integral, evaluated on EVQ's non-uniform cosh-quantile grid, naturally acquires a $1/\rho$ Jacobian factor, producing discrete $\chi^2$.

\paragraph{Closed form.} For $p{=}1$, the stiffness admits
\begin{equation}
S_{\chi^2}(\tau) = \frac{1}{d_{\mathrm{head}}}\!\left[\frac{\sinh\tau\cdot\arctan(\sinh\tau)}{\tau^2} - 1\right],
\end{equation}
verified numerically to relative error ${<}7{\times}10^{-6}$.

\paragraph{Exponent-matched diagnostic.} As a sensitivity analysis, one can ask which member of the $S_p$ family would match the target exponent $\tau_* \propto 1/\sqrt{L}$ exactly over this grid. The resulting $S_{\mathrm{derived}}(\tau)$ lies near $p \approx 0.85$ (high-resolution sweep; consistent with $0.80$ from the coarser grid), close to but distinct from the Pearson $\chi^2$ choice $p{=}1$. We do not use this as a derivation of the stiffness; it is reported as sensitivity analysis. The full sweep is reproduced by \texttt{scripts/analysis/verify\_stiffness\_and\_regime.py} (pure NumPy, ${\sim}60$s).

\subsection{Effective window \texorpdfstring{$L_{\mathrm{eff}}^{J}$}{L\_eff\^J} via softmax-Jacobian Rayleigh quotient}
\label{sec:leff-derivation}

\paragraph{From density perturbation to RoPE frequency shift.}
The EVQ pure-tether CDF is $F_\tau(\phi) = 1 - \sinh(\tau(1{-}\phi))/\sinh\tau$, with inverse-CDF warp $\Phi_\tau(u) = 1 - \tau^{-1}\operatorname{arsinh}((1{-}u)\sinh\tau)$. Setting $\theta{=}\tau^2$ and expanding $\rho_\tau(\phi) = 1 + \theta\, a(\phi) + O(\theta^2)$ with $a(\phi) = (1{-}\phi)^2/2 - 1/6$, the quantile displacement takes the closed form
\begin{equation}
\Phi_\tau(u) \;=\; u + \tau^2\, h(u) + O(\tau^4), \qquad h(u) \;=\; -\!\int_0^u a(v)\,dv \;=\; -\tfrac{1}{6}u(1{-}u)(2{-}u),
\label{eq:quantile-displacement}
\end{equation}
so RoPE frequencies perturb as $\omega_k(\tau) = \omega_k^{(0)}[1 - \tau^2(\log b)\,h(u_k) + O(\tau^4)]$ with $\omega_k^{(0)}{=}b^{-u_k}$. Hence $\dot\omega_k \coloneqq \partial_{\theta}\omega_k|_{\tau=0} = -(\log b)\,\omega_k^{(0)} h(u_k)$.

\paragraph{Explicit RoPE logit transport direction.}
Fix a layer/head and query $i$; write each $2$D RoPE pair as $q_{i,k}^{\mathbb{C}} = q_{i,2k{-}1}{+}iq_{i,2k}$, $k_{j,k}^{\mathbb{C}} = k_{j,2k{-}1}{+}ik_{j,2k}$, and $\alpha_{ij,k} = q_{i,k}^{\mathbb{C}}\overline{k_{j,k}^{\mathbb{C}}} = C_{ij,k} + iD_{ij,k}$. With relative position $r_{ij}$ and attention scale $s_{\mathrm{att}}$, the RoPE logit is $z_{ij}^{\mathrm{rope}}(\tau) = s_{\mathrm{att}}^{-1}\operatorname{Re}\sum_{k}\alpha_{ij,k}\exp(i r_{ij}\omega_k(\tau))$. The leading logit transport direction $g_i(j) \coloneqq \partial_\theta z_{ij}^{\mathrm{rope}}|_{\tau=0}$ has the closed form
\begin{equation}
g_i(j) \;=\; \frac{(\log b)\,r_{ij}}{s_{\mathrm{att}}}\sum_{k=1}^{K} \omega_k^{(0)}\,h(u_k)\!\left[\,C_{ij,k}\sin(r_{ij}\omega_k^{(0)}) + D_{ij,k}\cos(r_{ij}\omega_k^{(0)})\,\right].
\label{eq:g-closed-form}
\end{equation}
This is the Fr\'echet derivative of the logit with respect to the EVQ allocation perturbation at fixed activations. A finite-$\tau$ implementation that avoids first-order phase expansion at large $r_{ij}\omega$ is
\begin{equation}
g_i^{(\tau)}(j) \;=\; \frac{1}{s_{\mathrm{att}}\tau^2}\operatorname{Re}\sum_{k=1}^{K}\alpha_{ij,k}\!\left[\exp(ir_{ij}b^{-\Phi_\tau(u_k)}) - \exp(ir_{ij}b^{-u_k})\right],
\label{eq:g-finite-tau}
\end{equation}
with $g_i(j) = \lim_{\tau\to 0} g_i^{(\tau)}(j)$.

\paragraph{Primary definition: Rayleigh quotient of the softmax Jacobian.}
For the trained attention distribution $p_i(j)$ on keys $j\!\in\!\mathcal{K}_i$, $J_{\mathrm{sm}}(p_i) = \mathrm{diag}(p_i) - p_ip_i^\top$ satisfies
$g_i^\top J_{\mathrm{sm}}(p_i) g_i = \mathrm{Var}_{j\sim p_i}[g_i(j)] = \tfrac{1}{2}\sum_{j,k} p_i(j)p_i(k)(g_i(j)-g_i(k))^2.$
Since $J_{\mathrm{sm}}(p_i)\mathbf{1}{=}0$, we center $g_i$ by the key-wise mean operator $P_i$. The effective length is then defined by the \emph{ratio-of-means}
\begin{equation}
\frac{1}{L_{\mathrm{eff}}^{J}} \;=\; \kappa_{\mathrm{att}} \;=\; \frac{\mathbb{E}_{\ell,h,i}\!\left[\,g_i^\top J_{\mathrm{sm}}(p_i)g_i\,\right]}{\mathbb{E}_{\ell,h,i}\!\left[\,\|P_i g_i\|_2^2\,\right]},
\label{eq:leff-ratio-of-means}
\end{equation}
This is the directional curvature of the \emph{schedule-change} KL: for
$p_\epsilon=\operatorname{softmax}(z+\epsilon g)$,
$D_{\mathrm{KL}}(p\|p_\epsilon)=\tfrac12\epsilon^2g^\top
J_{\mathrm{sm}}(p)g+O(\epsilon^3)$, and for diffuse $p{=}1/L$ the quotient
reduces to $1/L$. The phase-variance utility $Q_1$ in
Proposition~\ref{prop:softmax-transport} is a distinct pairwise positional
functional. Replacing its $1/L$ factor by $1/L_{\mathrm{eff}}^J$ is therefore
a possible modelling extension, not a consequence of the KL identity and not
an empirical result used by the paper.

\subsection{Pure-tether branch: Fisher-forcing coefficient and residual bound}
\label{sec:pure-tether-bound}

\paragraph{Setup.} The full stationary ODE under $\mathcal{C}_{\mathrm{app}}$ plus a Fisher utility is $\rho''-\tau^2\rho = \gamma\,b^{-2\phi}$. Writing $s{=}2\log b$ and $q_s(\phi){=}e^{-s\phi}$, the pure-tether branch $\rho_{\mathrm{EVQ}}(\phi) = \tau\cosh(\tau(1{-}\phi))/\sinh\tau$ solves the homogeneous problem; the full solution is $\rho_{\mathrm{full}} = \rho_{\mathrm{EVQ}} + \delta\rho_{F}$.

\paragraph{Explicit form of $\gamma(b)$.} The Fisher utility in the Euler equation $\alpha\rho + \beta(G\rho) - \lambda_F V_b + \nu = 0$ has $V_b(\phi) = \eta_F(\phi)\,b^{-2\phi}$, where $\eta_F(\phi)$ is the activation-conditioned Fisher coefficient with
\begin{equation}
\eta_F(\phi_k) \;=\; \frac{1}{2\,s_{\mathrm{att}}^2}\,\mathbb{E}_{\ell,h,i,j}\!\left[\,w_{\ell h ij}\,|\alpha_{\ell h ij,k}|^2\,\right],\qquad \alpha_{\ell h ij,k} \;=\; q_{i,k}^{\mathbb{C}}\,\overline{k_{j,k}^{\mathbb{C}}}.
\label{eq:eta-F}
\end{equation}
Using $(G\rho)''{=}-\rho$, differentiating the Euler equation twice gives $\rho'' - \tau^2\rho = (\lambda_F/\alpha) V_b''$. For slowly varying $\eta_F$ compared with $e^{-s\phi}$,
\begin{equation}
\boxed{\;\gamma(b) \;=\; \frac{\lambda_F}{\alpha}\,\bar\eta_F\,s^2
\left[1+O\!\left(\tfrac{|\eta_F'|}{s\bar\eta_F}
+\tfrac{|\eta_F''|}{s^2\bar\eta_F}\right)\right],\qquad s=2\log b.\;}
\label{eq:gamma-explicit}
\end{equation}
The full $\eta_F(\phi)$-local version is $\gamma_{\mathrm{loc}}(\phi;b) = (\lambda_F/\alpha)[\eta_F''(\phi) - 2s\eta_F'(\phi) + s^2\eta_F(\phi)]$.

\paragraph{Suppression is L$^1$/CDF, not pointwise.} The particular solution $\delta\rho_{F}^{\mathrm{part}} = \gamma(b)\,q_s/(s^2-\tau^2)$ combined with~\eqref{eq:gamma-explicit} gives pointwise amplitude
\begin{equation}
\delta\rho_{F}^{\mathrm{part}}(\phi) \;\approx\; \frac{\lambda_F\bar\eta_F}{\alpha}\,\frac{s^2}{s^2-\tau^2}\,e^{-s\phi},
\label{eq:delta-rho-F-pointwise}
\end{equation}
so the $(s^2-\tau^2)^{-1}$ factor is cancelled by the $s^2$ in $\gamma(b)$, and $\delta\rho_{F}^{\mathrm{part}}$ is \emph{not} pointwise small in $s$. The true suppression comes from the boundary-layer width of $q_s$: $\|q_s\|_1 = (1-e^{-s})/s = O(1/\log b)$ and $\|q_s\|_2 = O(1/\sqrt{\log b})$. Let $C_{\mathrm{bc}}^{(p)}(\tau)$ denote the finite $L^p$ operator norm of the homogeneous correction enforcing the boundary and mass constraints, away from resonance. Then
\begin{align}
\|\rho_{\mathrm{full}}-\rho_{\mathrm{EVQ}}\|_1 \;&\le\; C_{\mathrm{bc}}^{(1)}(\tau)\;\frac{\lambda_F\bar\eta_F}{\alpha}\,\frac{s^2}{|s^2-\tau^2|}\,\frac{1-e^{-s}}{s} \;=\; O\!\left(\tfrac{1}{\log b}\right),\label{eq:tether-L1-correct}\\
\|\rho_{\mathrm{full}}-\rho_{\mathrm{EVQ}}\|_2 \;&\le\; C_{\mathrm{bc}}^{(2)}(\tau)\;\frac{\lambda_F\bar\eta_F}{\alpha}\,\frac{s^2}{|s^2-\tau^2|}\,\sqrt{\tfrac{1-e^{-2s}}{2s}} \;=\; O\!\left(\tfrac{1}{\sqrt{\log b}}\right),
\label{eq:tether-L2-correct}
\end{align}
and the CDF-level error is $\|F_{\mathrm{full}}-F_{\mathrm{EVQ}}\|_\infty \le \|\delta\rho_{F}\|_1 = O(1/\log b)$. For fixed Fisher coefficient and away from resonance, forcing is therefore small in CDF / $L^1$ transport at large $b$, but not necessarily pointwise small near $\phi{=}0$.

\paragraph{Warp amplification at large $\tau$.} Via the inverse-CDF map,
\begin{equation}
\|\phi_{\mathrm{full}}-\phi_{\mathrm{EVQ}}\|_\infty \;\le\; \frac{\|\delta\rho_{F}\|_1}{\inf_\phi \rho_{\mathrm{full}}(\phi)}, \qquad \inf_\phi \rho_{\mathrm{EVQ}}(\phi) \;=\; \frac{\tau}{\sinh\tau},
\label{eq:warp-bound}
\end{equation}
so, when the residual is small enough that $\rho_{\mathrm{full}}$ remains positive and comparable to $\rho_{\mathrm{EVQ}}$, CDF perturbations can be amplified at the scale $\sinh\tau/\tau$. The bound does not justify omitting the forcing branch at large $\tau$.

\paragraph{Dimensionless diagnostic.} Using $\|\rho_{\mathrm{EVQ}}-1\|_2 = \tau^2/\sqrt{45}+O(\tau^4)$ and $\Gamma_b \coloneqq (\lambda_F\bar\eta_F/\alpha)\,s^2/(|s^2-\tau^2|\tau^2)$, the $L^2$-level ratio is
\begin{equation}
R_F \;\coloneqq\; \frac{\|\rho_{\mathrm{full}}-\rho_{\mathrm{EVQ}}\|_2}{\|\rho_{\mathrm{EVQ}}-1\|_2} \;\le\; C_{\mathrm{bc}}^{(2)}\,\Gamma_b\,\sqrt{45}\,\sqrt{\tfrac{1-e^{-2s}}{2s}} \;=\; O\!\left(\tfrac{1}{\sqrt{\log b}}\right).
\label{eq:RF}
\end{equation}

\paragraph{Second-order utility deficit.} Let $A_\tau$ be the constrained
Hessian of the surrogate on zero-mass perturbations and let $q_s^\perp$ be the
projection of $q_s$ onto that tangent space. The full branch improves the
quadratic objective by
\begin{equation}
\Delta_F \;=\; \tfrac{1}{2}\gamma(b)^2\,\langle q_s^\perp,\,A_\tau^{-1} q_s^\perp\rangle,
\label{eq:tether-deficit}
\end{equation}
so dropping the forcing term is $O(\gamma(b)^2)$: if $R_F\!\le\!r$, the missed correction is $O(r^2)$ relative to the tether deformation energy.

\subsection{Channel-load motivation for Pearson \texorpdfstring{$\chi^2$}{chi2}}
\label{sec:chi2-load}

\paragraph{Normalization convention.} Throughout this paper, $S_{\chi^2}(\tau)$ denotes the $d_{\mathrm{head}}$-normalized stiffness $S_{\chi^2}(\tau){:=}\frac{1}{d_{\mathrm{head}}}\!\int_0^1\!(1{-}\rho_\tau)^2/\rho_\tau\,d\phi$, with leading expansion $\tau^4/(45\,d_{\mathrm{head}}){+}O(\tau^6)$ consistent with Proposition~\ref{prop:softmax-transport}. Identities below that produce a bare $\tau^4/45$ leading coefficient apply to $d_{\mathrm{head}}\,S_{\chi^2}(\tau)$ (the un-normalized integral); the $d_{\mathrm{head}}$ factor enters via the surrogate diagonal $\alpha{=}1/d_{\mathrm{rot}}$ ($={}1/d_{\mathrm{head}}$ for MHA via \S\ref{sec:tau-scaling}; ${=}1/d_{\mathrm{rot}}$ for compressed-RoPE MLA via \S\ref{sec:mla-results}).

\paragraph{Channel-load axioms.} Inverse-CDF quantization places channels with cell width $\ell(\phi){=}1/\rho(\phi)$ and samples channels according to $dP_\rho{=}\rho\,d\phi$. Under the modeling axioms (i)~channel-measure averaging, (ii)~load $\ell$ as the primitive local quantity, (iii)~second central moment, (iv)~uniform density has zero stiffness, (v)~local additivity, the canonical load-variance penalty is
\begin{equation}
S(\rho) \;=\; \mathrm{Var}_{\phi\sim\rho\,d\phi}\!\left[\,\tfrac{1}{\rho(\phi)}\,\right] \;=\; \int_0^1 \frac{(1-\rho(\phi))^2}{\rho(\phi)}\,d\phi \;=\; S_{\chi^2}(\rho).
\label{eq:chi2-load-identity}
\end{equation}
We treat this as a canonical motivated choice under the stated axioms, not a uniqueness claim.

\paragraph{Load-moment exponent.} Distinct from the weighted-$L^2$ family $S_p(\tau){=}\frac{1}{d_{\mathrm{head}}}\int(\rho_\tau{-}1)^2/\rho_\tau^p\,d\phi$ used in Table~\ref{tab:stiffness-sweep} (where Pearson $\chi^2$ is $p{=}1$), we also examine a load-moment family $S^{\mathrm{mom}}_p(\rho){=}\mathbb{E}_\rho[(1/\rho{-}1)^p]$. Substituting the EVQ cosh density and minimizing $S^{\mathrm{mom}}_p{-}\lambda U$ gives $\tau^{2(p-1)}\asymp d_{\mathrm{head}}^2/L_{\mathrm{eff}}^{J}$, hence $\tau^{*}\propto (L_{\mathrm{eff}}^{J})^{-1/(2p-2)}$; only $p{=}2$ yields the observed $L^{-1/2}$ exponent within this family. The two families parametrize different $f$-divergence neighborhoods and need not match on the finite-$\tau$ grid.

\paragraph{$f$-divergences at finite $\tau$.} Smooth $f$-divergences (KL, Jeffreys, R\'enyi) are locally quadratic around $\rho{=}1$. Over the finite grid used in Table~\ref{tab:stiffness-sweep}, however, their fitted exponents range from $0.58$ to $0.81$. This is model sensitivity, not a trained-PPL comparison.

\subsection{Discrete-channel transport gap and the MLA channel-scarcity mechanism}
\label{sec:discrete-continuous-gap}

\paragraph{Setup.} Inverse-CDF quantization replaces the continuous density $\rho$ by an atomic measure. Let $F(\phi){=}\int_0^\phi\rho$, $Q(u){=}F^{-1}(u)$, midpoint grid $u_k{=}(k{-}\tfrac{1}{2})/K$, and $\phi_k{=}Q(u_k)$. The empirical channel measure is $\mu_K{=}K^{-1}\sum_{k=1}^{K}\delta_{\phi_k}$. Direct $\|\rho_K-\rho\|$ is ill-posed: one cannot subtract an atomic measure from a density. We therefore work in Wasserstein distance for transport and in a quantile-cell histogram $\rho_K$ for density.

\paragraph{Transport bounds.} Assume $0<m\le\rho\le M$ and $\|\rho'\|_\infty\le B$; then $Q$ is $(1/m)$-Lipschitz. The midpoint quantization errors satisfy
\begin{equation}
W_\infty(\mu_K,\mu_\rho) \;\le\; \frac{1}{2Km}, \qquad W_1(\mu_K,\mu_\rho) \;\le\; \frac{1}{4Km}.
\label{eq:wasserstein-bound}
\end{equation}
Defining the quantile-cell histogram $\rho_K(\phi){=}1/(K|I_k|)$ on $I_k{=}[Q((k{-}1)/K),\,Q(k/K)]$, the histogram density satisfies
\begin{equation}
\|\rho_K-\rho\|_1 \;\le\; \frac{B}{Km}, \qquad \|\rho_K-\rho\|_\infty \;\le\; \frac{B}{Km}.
\label{eq:density-bound}
\end{equation}

\paragraph{Application to EVQ-Cosh.} For $\rho_\tau(\phi)=\tau\cosh(\tau(1{-}\phi))/\sinh\tau$, $m_\tau=\tau/\sinh\tau$ and $B_\tau=\tau^2$. Substituting,
\begin{equation}
W_1(\mu_{K,\tau},\mu_{\rho_\tau}) \;\le\; \frac{\sinh\tau}{4K\tau}, \qquad \|\rho_{K,\tau}-\rho_\tau\|_1 \;\le\; \frac{\tau\sinh\tau}{K} \;=\; \frac{\tau^2}{K}+O\!\left(\tfrac{\tau^4}{K}\right).
\label{eq:evq-discrete-bound}
\end{equation}

\paragraph{Kernel-integrated errors.} For a smooth two-variable kernel $K_{\mathrm{sm}}$ with Lipschitz constant $L_K$ in each argument,
\begin{equation}
\left|\iint K_{\mathrm{sm}}\,d\mu_K\,d\mu_K - \iint K_{\mathrm{sm}}\,d\mu_\rho\,d\mu_\rho\right| \;\le\; \frac{L_K}{2Km} + O(K^{-2}).
\label{eq:kernel-bound}
\end{equation}
The singular $\delta$-component of $K_{\mathrm{app}}$ must be interpreted as a cell average (not as a literal atom-kernel pairing), after which the bound applies to the smoothed kernel.

\paragraph{High-resolution distortion.} For a smooth weighted distortion $w$,
\begin{equation}
\mathcal{D}_K[\rho] \;\approx\; \frac{1}{12K^2}\int_0^1 \frac{w(\phi)}{\rho(\phi)^2}\,d\phi + O(K^{-3}).
\label{eq:high-res-distortion}
\end{equation}
This is the finite-channel mechanism: a shaped $\rho$ can reduce the weighted inverse-density load $\int w/\rho^2$; the absolute benefit scales as $K^{-2}$ for smooth quantization error and as $K^{-1}$ for transport/kernel terms.

\paragraph{Channel scarcity amplifies the density-shape correction.} The
standard MHA configuration has $d_{\mathrm{head}}{=}d_{\mathrm{rope}}{=}64$
and $K{=}32$; the MLA study has $d_{\mathrm{head}}{=}64$,
$d_{\mathrm{rope}}{=}32$, $d_{\mathrm{nope}}{=}32$ and $K{=}16$. Halving
$K$ scales the $K^{-2}$ component of the distortion bound by $4\times$ and
the $K^{-1}$ component by $2\times$. The observed MLA gain
(\S\ref{sec:mla-results}) is a test of the predicted direction, not a controlled
quantitative test of either bound.

\paragraph{Caveat.} At large $\tau$, $m_\tau{=}\tau/\sinh\tau$ becomes small and the $1/(Km_\tau)$ factor grows. The transport bounds can therefore become vacuous in the same regime where the Fisher-residual analysis no longer guarantees a small warp error (\S\ref{sec:pure-tether-bound}).
